\documentclass[10pt,a4paper]{amsart}
\usepackage[margin=19mm]{geometry}
\usepackage{amsmath,amssymb,mathtools}
\usepackage[scr=rsfs,cal=euler]{mathalfa}
\usepackage{xfrac}
\usepackage[unicode,hidelinks,bookmarks=false]{hyperref}
\newcommand{\ie}{i.e.\ }
\newcommand{\cf}{cf.\ }
\newcommand{\resp}{respectively\ }
\newcommand{\Subsection}{\subsection}
\providecommand{\nobreakdash}{\nobreak\hbox{-}}
\setlength{\emergencystretch}{2em}
\multlinegap0pt
\usepackage{bm}
\usepackage{bbm}
\usepackage{dsfont}
\usepackage{xspace}
\usepackage{cancel}
\usepackage{mathtools}
\usepackage[shortlabels]{enumitem}
\usepackage{amsthm}
\makeatletter
\@ifundefined{theorem}{\newtheorem{theorem}{Theorem}}{}
\@ifundefined{definition}{\newtheorem{definition}[theorem]{Definition}}{}
\@ifundefined{proposition}{\newtheorem{proposition}[theorem]{Proposition}}{}
\makeatother
\newcommand{\Ga}{\Gamma}
\newcommand{\N}{\mathbb N}
\newcommand{\Om}{\Omega}
\newcommand{\abs}[1]{|#1|}
\newcommand{\dG}{\;\mathrm d\Ga}
\newcommand{\free}{\mathrm{free}}
\newcommand{\intG}{\int_\Gamma}
\title[On a bulk-surface Navier-Stokes-Cahn-Hilliard model: Existen]{On a bulk-surface Navier-Stokes-Cahn-Hilliard model: Existence of weak solutions and asymptotic limits}
\author{Jonas Stange}
\date{Source: arXiv:2608.24647v1 (CC BY 4.0)}
\begin{document}
\maketitle

\noindent\emph{Source and licence.} On a bulk-surface Navier-Stokes-Cahn-Hilliard model: Existence of weak solutions and asymptotic limits. Version arXiv:2608.24647v1 (CC BY 4.0), distributed under the Creative Commons Attribution 4.0 International licence, as recorded in the arXiv licence metadata for this exact version and shown on the abstract page https://arxiv.org/abs/2608.24647v1. Full rights notice: licenses/arxiv-2608.24647-CC-BY-4.0.txt.\par\smallskip

\noindent\textit{Editorial scope.} Counted unit: the Proposition whose source label is \texttt{Proposition:MoscoConvergence} in the article (source0.tex lines 996--998, source label \texttt{Proposition:MoscoConvergence}), delivered together with the article's own complete original proof of it (source0.tex lines 1000--1070, one \texttt{proof} environment of 71 lines). No mathematics is written, repaired or re-derived: statement and proof are the article's own text line for line. Required context retained uncounted: Mosco1 (lines 979--991); Mosco2 (lines 983--985). Every definition, display and label of those lines is retained unchanged. Omitted from this package and not redistributed: 1--978, 986--995, 999--999, 1071--2333. Nothing inside a retained excerpt is omitted or altered: every retained body is delivered exactly as the article's lines are written, so no omission receipt of any kind accompanies this package. Citations: the retained excerpts carry no citation command at all, so no bibliography entry of the article is transcribed and no reference is redistributed. Reference adaptation: none was needed and none was made; every reference inside the delivered unit resolves inside it, so no unresolved reference or citation is produced. Printed slips kept literal: no printed word, symbol, hypothesis or number of the counted statement or of its proof was altered. Layout and class: the article's own class is \texttt{amsart} with packages including amsfonts, eurosym, amssymb, amsthm, amsmath, amsaddr, bm, cite, mathrsfs, xcolor, fontenc, geometry, hyperref, mathtools; the wrapper is the lane's pinned \texttt{amsart} editorial wrapper at 10pt on A4, which restates the theorem-like environments the retained text needs and adds the article's own macro definitions that the retained text needs (\texttt{\textbackslash Ga}, \texttt{\textbackslash N}, \texttt{\textbackslash Om}, \texttt{\textbackslash abs}, \texttt{\textbackslash dG}, \texttt{\textbackslash free}, \texttt{\textbackslash intG}), with extra wrapper packages (bm, bbm, dsfont, xspace, cancel, mathtools, enumitem). The delivered numbering is the unit's own. Assets: no retained span contains a figure environment, a graphics-inclusion command, a PDF-page inclusion, a table or a TikZ picture; no image file is copied into this package, the package declares no asset dependency and no image file is redistributed with it. Licence: version arXiv:2608.24647v1 of this article is distributed under the Creative Commons Attribution 4.0 International licence, as recorded in the arXiv licence metadata for this exact version and shown on the abstract page https://arxiv.org/abs/2608.24647v1. A full rights notice is delivered at licenses/arxiv-2608.24647-CC-BY-4.0.txt. Characters: every retained excerpt is pure ASCII, so the delivered engine, XeLaTeX with its default Latin Modern text font, reports no missing character.
\begin{definition}
    Let $H$ be a Hilbert space, and let $\Phi,\Phi_n:H\rightarrow [0,+\infty]$, $n\in\N$, be proper, lower semicontinuous and convex functionals on $H$. We say that $(\Phi_n)_{n\in\N}$ converges to $\Phi$ in the sense of Mosco if the following two conditions hold:
    \begin{enumerate}[label=\textnormal{\bfseries(M\arabic*)}]
        \item \label{Mosco1} For every sequence $(x_n)_{n\in\N}\subset H$ converging weakly to $x\in H$ it holds that
        \begin{align*}
            \Phi(x) \leq \liminf_{n\rightarrow\infty} \Phi_n(x_n).
        \end{align*}
        \item \label{Mosco2} For every $x\in H$ there exists a sequence $(x_n)_{n\in\N}\subset H$ converging strongly to $x$ such that
        \begin{align*}
            \Phi(x) \geq \limsup_{n\rightarrow\infty} \Phi_n(x_n).
        \end{align*}
    \end{enumerate}
\end{definition}

\begin{proposition}\label{Proposition:MoscoConvergence}
    The functional $\widetilde E_\free^K$ converges in the sense of Mosco to $\widetilde E_\free^0$ and $\widetilde E_\free^\infty$, respectively, as $K\rightarrow 0$ and $K\rightarrow\infty$, respectively.
\end{proposition}

\begin{proof}
    We first consider the case $K\rightarrow 0$. To this end, let $\{K_n\}_{n\in\N}\subset(0,\infty)$ be a sequence with $K_n\rightarrow 0$ as $n\rightarrow\infty$, and let $\{(u_n,v_n)\}_{n\in\N}\subset\mathcal{L}^2$ be a sequence converging weakly in $\mathcal{L}^2$ to $(u,v)\in\mathcal{L}^2$ as $n\rightarrow\infty$. Without loss of generality, we can assume that 
    \begin{align*}
        \liminf_{n\rightarrow\infty} \widetilde E_\free^{K_n}(u_n,v_n) < + \infty,
    \end{align*}
    as otherwise \ref{Mosco1} is trivially fulfilled. Then, going over to a subsequence, we have 
    \begin{align*}
        \lim_{k\rightarrow\infty} \widetilde E_\free^{K_{n_k}}(u_{n_k},v_{n_k}) = \liminf_{n\rightarrow\infty} \widetilde E_\free^{K_n}(u_n,v_n),
    \end{align*}
    as well as
    \begin{subequations}
        \begin{alignat}{2}
            (u_{n_k},v_{n_k}) &\rightarrow (u,v) &&\qquad\text{weakly in~}\mathcal{H}^1, \label{Prelim:Mosco:WeakConv} \\
            (u_{n_k},v_{n_k}) &\rightarrow (u,v) &&\qquad \text{strongly in~}\mathcal{L}^2 \label{Prelim:Mosco:StrongConv}
        \end{alignat}
    \end{subequations}
    as $k\rightarrow\infty$. Moreover, as the energy is bounded along this subsequence, it in particular holds that
    \begin{align}
        \sup_{k\in\N} \frac{1}{K_{n_k}}\intG \frac12(\alpha v_{n_k} - u_{n_k})^2\dG < +\infty.
    \end{align}
    Consequently, we readily infer that
    \begin{align*}
        \alpha v_{n_k} - u_{n_k} \rightarrow 0 \qquad\text{strongly in~}L^2(\Ga)
    \end{align*}
    as $k\rightarrow\infty$. Since we already know that \eqref{Prelim:Mosco:WeakConv} holds, the continuity of the trace operator with the above convergence entails that $u = \alpha v$ a.e. on $\Ga$. Additionally, by possibly going over to another subsequence, we readily infer from $\abs{u_{n_k}}\leq 1$ a.e.~in $\Om$ and $\abs{v_{n_k}}\leq 1$ a.e.~on $\Ga$ that $\abs{u}\leq 1$ a.e.~in $\Om$ and $\abs{v}\leq 1$ a.e.~on $\Ga$, respectively. Then, from the strong convergence \eqref{Prelim:Mosco:StrongConv} and the uniform continuity of $F_0$ and $G_0$, the continuity of the Nemytskii operator directly gives that
    \begin{subequations}
        \begin{alignat}{2}
            F_0(u_{n_k})&\rightarrow F_0(u) &&\qquad\text{strongly in~}L^1(\Omega), \label{Prelim:Mosco:Conv:Pot:F:K0}\\
            G_0(v_{n_k})&\rightarrow G_0(v) &&\qquad\text{strongly in~}L^1(\Ga)\label{Prelim:Mosco:Conv:Pot:G:K0}
        \end{alignat}
    \end{subequations}
    as $k\rightarrow\infty$. Thus, we deduce with \eqref{Prelim:Mosco:WeakConv} and \eqref{Prelim:Mosco:Conv:Pot:F:K0}-\eqref{Prelim:Mosco:Conv:Pot:G:K0}, and the weak lower semicontinuity of norms that
    \begin{align*}
        \widetilde E_\free^0(u,v) \leq \liminf_{k\rightarrow\infty} \widetilde E_\free^{K_{n_k}}(u_{n_k},v_{n_k}) = \lim_{k\rightarrow\infty} \widetilde E_\free^{K_{n_k}}(u_{n_k},v_{n_k}) = \liminf_{n\rightarrow\infty} \widetilde E_\free^{K_n}(u_n,v_n).
    \end{align*}
    
    For \ref{Mosco2}, let $(u,v)\in\mathcal{L}^2$. Without loss of generality, we can assume that $(u,v)\in\mathrm{Dom}\,\widetilde E_\free^0$, as otherwise $\widetilde E_\free^0(u,v) = +\infty$. We can thus choose the constant recovery sequence $(u_n,v_n)\coloneqq (u,v)\in \mathrm{Dom}\,\widetilde E_\free^0\subset \mathrm{Dom}\,\widetilde E_\free^{K_n}$ for all $n\in\N$, and readily obtain
    \begin{align*}
        \widetilde E_\free^0(u,v) = \limsup_{n\rightarrow\infty} \widetilde E_\free^{K_n}(u,v) = \limsup_{n\rightarrow\infty} \widetilde E_\free^{K_n}(u_n,v_n),
    \end{align*}
    which proves the Mosco convergence of $\widetilde E_\free^K$ as $K\rightarrow 0$ and finishes the proof in this case.

    Now, consider $K_n\rightarrow\infty$ as $n\rightarrow\infty$. Let $\{(u_n,v_n)\}_{n\in\N}\subset\mathcal{L}^2$ be a sequence converging weakly in $\mathcal{L}^2$ to $(u,v)\in\mathcal{L}^2$ as $n\rightarrow\infty$. Without loss of generality we can assume that 
    \begin{align*}
        \liminf_{n\rightarrow\infty} \widetilde E_\free^{K_n}(u_n,v_n) < + \infty,
    \end{align*}
    as otherwise \ref{Mosco1} is trivially fulfilled. Then, as for $K\rightarrow 0$, we find a subsequence such that
    \begin{align*}
        \lim_{k\rightarrow\infty} \widetilde E_\free^{K_{n_k}}(u_{n_k},v_{n_k}) = \liminf_{n\rightarrow\infty} \widetilde E_\free^{K_n}(u_n,v_n),
    \end{align*}
    and
    \begin{subequations}
        \begin{alignat}{2}
            (u_{n_k},v_{n_k}) &\rightarrow (u,v) &&\qquad\text{weakly in~}\mathcal{H}^1, \label{Prelim:Mosco:WeakConv:infty} \\
            (u_{n_k},v_{n_k}) &\rightarrow (u,v) &&\qquad \text{strongly in~}\mathcal{L}^2 \label{Prelim:Mosco:StrongConv:infty}
        \end{alignat}
    \end{subequations}
    as $k\rightarrow\infty$. Then, the same Nemytskii argument and weak lower semicontinuity give
    \begin{align*}
        \widetilde E_\free^\infty(u,v) \leq \liminf_{k\rightarrow\infty} \widetilde E_\free^\infty(u_{n_k},v_{n_k}) \leq \liminf_{k\rightarrow\infty} \widetilde E_\free^{K_{n_k}}(u_{n_k},v_{n_k}) = \liminf_{n\rightarrow\infty} \widetilde E_\free^{K_n}(u_n,v_n).
    \end{align*}
    For \ref{Mosco2}, let $(u,v)\in\mathcal{L}^2$. Without loss of generality we can again assume that $(u,v)\in\mathrm{Dom}\,\widetilde E_\free^\infty$, as otherwise $\widetilde E_\free^\infty(u,v) = +\infty$. We then choose again the constant sequence $(u_n,v_n)\coloneqq (u,v) \in\mathrm{Dom}\,\widetilde E_\free^\infty = \mathrm{Dom}\,\widetilde E_\free^{K_n}$ for all $n\in\N$, and observe that
    \begin{align*}
        \frac{1}{K_n}\intG \frac12(\alpha v_n - u_n)^2\dG = \frac{1}{K_n}\intG \frac12(\alpha v - u)^2\dG \rightarrow 0
    \end{align*}
    as $n\rightarrow\infty$, which entails that
    \begin{align*}
        \widetilde E_\free^\infty(u,v) = \limsup_{n\rightarrow\infty} \widetilde E_\free^{K_n}(u,v) = \limsup_{n\rightarrow\infty} \widetilde E_\free^{K_n}(u_n,v_n).
    \end{align*}
    This shows the Mosco convergence of $\widetilde E_\free^K$ as $K\rightarrow\infty$ and finishes the proof.
\end{proof}

\end{document}
